\section{Related Work}
\textbf{Continual Learning}
Various continual learning methods have been proposed to mitigate the problem of CF. The methods can be broadly divided into architecture-based methods \cite{yoon2018lifelong,serra18a}, regularization-based methods \citet{li2017learning,kirkpatrick2017overcoming}, and rehearsal-based methods \cite{NEURIPS2019_fa7cdfad}.
Under the fixed model structure, regularization-based methods \cite{kirkpatrick2017overcoming,aljundi2018memory,li2017learning} optimize network parameters on the current task while constraining the representation drift. For example, \citet{li2017learning} propose learning without forgetting (LwF) to tackle this problem, which regularizes the model output of current data close to those trained for the previous model. Another category of fixed-structure strategies (rehearsal-based) stores a limited subset of samples from previous tasks to mitigate CF such as ER \cite{NEURIPS2019_fa7cdfad}, RM \cite{Bang_2021_CVPR}, and iCaRL \cite{Rebuffi_2017_CVPR}.

% Architecture-based models rely on the dynamic model structure to solve new coming tasks, i.e. allocate different modules for different training tasks \cite{yoon2018lifelong,serra18a}. But the size of the model becomes increasingly large with the new learning tasks, and  knowledge cannot be transferred during the training.  Architecture-based models rely on dynamic model structure to solve new coming tasks, i.e. allocate different modules for different training tasks \cite{yoon2018lifelong,serra18a}. 


% For example, an adapter is inserted into the gap between transformer modules of pre-trained language models for a new coming task, and the parameters of the transformers are fixed to preserve the knowledge of pre-trained models.  Thus, the 
%  size of the model becomes increasingly large with the increasing number of learning tasks, and meanwhile, the  knowledge can not be transferred during the training. 

% The other two classes of continual learning methods are under the fixed model structure. Regularization-based methods \cite{kirkpatrick2017overcoming,aljundi2018memory,li2017learning} attempt to optimize network parameters on the current task while constraining the parameters drift of the previous model by adding some regularization terms. For example, \citet{li2017learning} proposes a method called learning without forgetting (LwF) to tackle this problem, which regularizes the model output of current data close to those trained for the previous model. Another category of continual learning (rehearsal-based methods) stores a limited subset of samples from previous tasks in the fixed-size memory buffer to mitigate CF such as ER \cite{NEURIPS2019_fa7cdfad}, RM \cite{Bang_2021_CVPR}, iCaRL \cite{Rebuffi_2017_CVPR}, GEM \cite{NIPS2017_f8752278} , AGEM \cite{chaudhry2018efficient}. It has shown promising performance for various downstream tasks of natural language processing \cite{NEURIPS2020_3fe78a8a,sun2019lamol,holla2020meta,wang-etal-2019-sentence}.



% It has shown promising performance for various downstream tasks of natural language processing \cite{NEURIPS2020_3fe78a8a,sun2019lamol,holla2020meta,wang-etal-2019-sentence}.


 % \citet{yu2020semantic} propose an embedding network to solve the class-incremental continual learning task. They adopt the training objective of triplet loss to force negative samples further apart than the positive ones from the anchor. And they use prototypes as the classification criterion, and they further propose a semantic drift compensation module to mitigate the representation drift. However, the model focus on the task of class-incremental continual learning, and semantic drift compensation relies on the assumption that the directions of prototypes' drift are not disjointed.

\textbf{Contrastive Learning}
Contrastive learning is initially introduced in self-supervised settings  and proved to  subsume or significantly outperform traditional contrastive losses such as triplet loss \cite{chen2020simple,wu2018unsupervised,gao-etal-2021-simcse,}.  For example, \citet{khosla2020supervised} first propose the idea of self-supervised contrastive learning and prove that the method is more robust to natural corruptions, stable to hyper-parameter settings,  and has strong transfer performance. \citet{luo2022mere} uses supervised contrastive learning combined with a kNN inference module for cross-domain sentiment analysis,  showing a stronger generalization ability compared with standard CE-based methods. 

\citet{cha2021co2l} propose a contrastive continual learning method, Co$^2$L, for class-incremental continual learning. The method uses an asymmetric supervised contrastive loss to enlarge the distance between representations of previous and new tasks.  However, there are significant differences between our model and Co$^2$L. First, the asymmetric contrastive loss of Co$^2$L is unsuitable for task-incremental continual learning,  because a representation can be predicted as different labels according to task objectives. Second, Co$^2$L  uses a decoupled classification layer for inference, i.e. it learns representations first and then learns a linear classifier separately, causing low extensibility and high complexity of the model. In contrast, we use a kNN module as the classification criteria to enhance the extensibility of the model. Third, Co$^2$L only considers the representation drift from the view of regularization. But we also mitigate the problem of representation drift by feeding memory data into the current model to obtain updated classification criteria.
